% Slide — DQN antes/depois
\begin{frame}{DQN -- \antes\ vs. \depois}
  \begin{columns}[T,onlytextwidth]
    \column{0.48\textwidth}
      \antes\\[0.3em]
      \small Rede \texttt{\_NN}: NumPy escrito à mão
      \begin{itemize}\footnotesize
        \item MLP simples, \textit{forward}/\textit{backward} manuais,
          \texttt{ReLU}, sem otimizador adaptativo (SGD puro).
        \item DQN padrão: mesma rede escolhe e avalia a ação do próximo
          estado -- \textbf{superestimação} de valor conhecida na
          literatura.
        \item Perda: erro quadrático médio.
        \item Sem \textit{clipping} de gradiente.
      \end{itemize}
    \column{0.48\textwidth}
      \depois\\[0.3em]
      \small \texttt{DoubleDQNAgent} (PyTorch)
      \begin{itemize}\footnotesize
        \item Arquitetura \highlight{\textit{dueling}} (Wang et al., 2016):
          $Q(s,a)=V(s)+[A(s,a)-\overline{A(s,\cdot)}]$.
        \item \highlight{Double DQN} (van Hasselt et al., 2016): rede
          \textit{online} escolhe a ação, rede \textit{alvo} avalia --
          $y=r+\gamma\,Q_{\text{alvo}}(s',\arg\max_a Q_{\text{online}}(s',a))$.
        \item Perda de \textbf{Huber} (\texttt{smooth\_l1\_loss}), Adam,
          \textit{grad clipping} em 10.
      \end{itemize}
  \end{columns}
  \vspace{0.6em}

  \begin{block}{Por que Double DQN dueling especificamente}
    \small Em horizontes curtos ($T=38$ ciclos), a superestimação do DQN
    clássico é especialmente danosa: o agente aprende a pedir demais
    porque valoriza excessivamente estados de estoque alto. A separação
    $V(s)/A(s,a)$ estabiliza o treino em estados onde a ação pouco altera
    o retorno -- frequente quando o estoque já está muito acima do ponto
    de pedido.
  \end{block}

  \note{
    O que era rotulado "DQN" em julho era, tecnicamente, DQN vanilla com
    uma rede neural implementada do zero em NumPy -- funcional, mas sem
    nenhuma das melhorias que a literatura de RL desenvolveu desde 2015
    especificamente para o problema de superestimação de valor. A versão
    atual não é apenas uma reescrita em PyTorch: é a adoção de Double DQN
    com arquitetura dueling, que são, juntos, considerados a versão de
    estado da arte razoável para este slot no portfólio -- van Hasselt et
    al. 2016 para o Double DQN, Wang et al. 2016 para dueling. A perda de
    Huber, em vez de erro quadrático, também é mais robusta a
    recompensas de escala variável entre séries.
  }
\end{frame}
